\documentclass[11pt,a4paper]{report}

% Essential packages
\usepackage{amsmath, amssymb, amsthm} % AMS Packages
\usepackage{graphicx,color}           % Packages for graphics and color
\usepackage[left=1.5in, right=1in, top=1in, bottom=1in, includefoot, headheight=13.6pt]{geometry}
\usepackage[square, comma, numbers, sort&compress]{natbib}

% Optional customization packages
\usepackage{lmodern}                  % Custom fonts
\usepackage[T1]{fontenc}              % Ensure correct font encoding

\usepackage{url}

\usepackage[hang, small, bf, margin=0pt, tableposition=bottom]{caption}
\setlength{\abovecaptionskip}{10pt}   % Custom captions
\graphicspath{{./img/}}

% Tables
\usepackage{xcolor}
\usepackage{colortbl}

\newcommand{\mc}[2]{\multicolumn{#1}{c}{#2}}
\definecolor{Gray}{gray}{0.83}

\newcolumntype{a}{>{\columncolor{Gray}}c}
\newcolumntype{b}{>{\columncolor{white}}c}

% Page layout
\parindent 0pt
\parskip 1ex
\renewcommand{\baselinestretch}{1.49}
\numberwithin{equation}{section}
\renewcommand{\bibname}{References}
\renewcommand{\contentsname}{Contents}
\pagenumbering{roman}

\bibliographystyle{unsrtnat}

\newcommand{\acrolabel}[1]{\makebox[3cm][l]{\textbf{#1}}}
\newenvironment{acronyms}{\begin{list}{}{\renewcommand{\makelabel}{\acrolabel}}}{\end{list}}

% \includeonly{tex/chapter1}          % Option to generate specific chapters

% Customising headers - fancyhdr.pdf
\usepackage{fancyhdr}
\pagestyle{fancy}
\rhead{}
\lhead{\nouppercase{\textsc{\leftmark}}}
\renewcommand{\headrulewidth}{0pt}
\makeatletter
\renewcommand{\chaptermark}[1]{\markboth{\textsc{\@chapapp}\ \thechapter:\ #1}{}}
\makeatother

% Customising chapter headings - sectsty.pdf
\usepackage{sectsty}
\chapterfont{\Large\sc\centering}
\chaptertitlefont{\centering}
\subsubsectionfont{\centering}

% Hyperreferencing and citations
\usepackage{hyperref}
\hypersetup{
    colorlinks,
    citecolor=blue,
    filecolor=blue,
    linkcolor=blue,
    urlcolor=blue
}

% Section symbol
\usepackage{cleveref}
\crefname{section}{\S}{\S\S}
\Crefname{section}{\S}{\S\S}
\crefname{subsection}{\S}{\S\S}
\Crefname{subsection}{\S}{\S\S}

% Actual document
\begin{document}

    % Titlepage
    \newpage
    \pdfbookmark[0]{Main Title}{maintitle}
    \begin{titlepage}
        \centering
        \vspace {1cm}
        \huge{\textbf{An Embodied Robotic Agent Learning Affordances with Intrinsic Motivations and Solving Extrinsic Tasks with Developmental Learning}} \\ [0.5cm]
        \begin{figure}[h]
\centering
\includegraphics[scale =0.9]{img/iub.png}
\end{figure}
        \vspace {0.5cm}
        \Large{By} \\
        \Large{\textbf{Fawwad Hassan Jaskani}} \\
        \Large{Reg No:} \\[0.75cm]
        \Large{Supervisor} \\
        \Large{\textbf{Dr. Nadia Rashid}} \\[0.75cm]
        \Large{Department of Computer Systems Engineering \\
        Faculty of Engineering \\
        Islamia University of Bahawalpur \\
        Pakistan} \\ [0.75 cm]
        \Large{September 2020}
    \end{titlepage}

    \newpage
    \pdfbookmark[0]{Title Page}{titlepage}
    \begin{titlepage}
        \centering
        \huge{\textbf{An Embodied Robotic Agent Learning Affordances with Intrinsic Motivations and Solving Extrinsic Tasks with Developmental Learning}} \\ [0.1cm]
      
        \Large{By} \\
        \Large{\textbf{Fawwad Hassan Jaskani}} \\
        \Large{Reg No:} \\[0.2cm]
        \Large{Supervisor} \\
        \Large{\textbf{Dr. Nadia Rashid}} \\
        \line(1,0){150} \\
        \Large{A thesis submitted in conformity with the requirements for \\
        the degree of \emph{Master of Science} in \\
       Computer Systems Engineering} \\[0.2cm]
        \Large{Department of Computer Systems Engineering \\
        Faculty of Engineering \\
        Islamia University of Bahawalpur \\
       Pakistan} \\
        \Large{September 2020}
    \end{titlepage}

    \newpage
    \pdfbookmark[0]{Declaration}{declaration}
    \chapter*{Declaration} % (fold)
    \label{cha:declaration}
    I, \textit{Fawwad Hassan Jaskani} declare that this thesis titled ``An Embodied Robotic Agent Learning Affordances with Intrinsic Motivations and Solving Extrinsic Tasks with Developmental Learning'' and the work presented in it are my own and has been generated by me as a result of my own original research. \\

    I confirm that:
    \begin{enumerate}
        \item This work was done wholly or mainly while in candidature for a Master of Science degree at Islamia University of Bahawalpur (IUB)
        \item Where any part of this thesis has previously been submitted for a degree or any other qualification at IUB or any other institution, this has been clearly stated
        \item Where I have consulted the published work of others, this is always clearly attributed
        \item Where I have quoted from the work of others, the source is always given. With the exception of such quotations, this thesis is entirely my own work
        \item I have acknowledged all main sources of help
        \item Where the thesis is based on work done by myself jointly with others, I have made clear exactly what was done by others and what I have contributed myself
    \end{enumerate}

    \begin{flushright}
        \line(1,0){100} \\
        Fawwad Hassan Jaskani, \\
        Reg No:
    \end{flushright}

    \newpage
    \pdfbookmark[0]{Copyright Notice}{copyright}
    \chapter*{Copyright Notice} % (fold)
    \label{cha:copyright_notice}
    \begin{itemize}
        \item Copyright in text of this thesis rests with the student author. Copies (by any process) either in full, or of extracts, may be made only in accordance with instructions given by the author and lodged in the Library of FoE, IUB. Details may be obtained by the Librarian. This page must form part of any such copies made. Further copies (by any process) may not be made without the permission (in writing) of the author.
        \item The ownership of any intellectual property rights which may be described in this thesis is vested in FoE, IUB, subject to any prior agreement to the contrary, and may not be made available for use by third parties without the written permission of FoE, which will prescribe the terms and conditions of any such agreement.
        \item Further information on the conditions under which disclosures and exploitation may take place is available from the Library of FoE, IUB, Bahawalpur.
    \end{itemize}
    % chapter copyright_notice (end)

    % Dedication
    \newpage
    \vspace*{8cm}
    \pdfbookmark[0]{Dedication}{dedication}
    \begin{center}
        \Large This thesis is dedicated to \emph{my beloved parents}
    \end{center}

    % Abstract
    \newpage
    \pdfbookmark[0]{Abstract}{abstract}
    \chapter*{Abstract}
Robot has gained an acceptable entity in modern world because it integrates sensors, actuation and intelligence phenomena’s. The modern world robot not only perform in industry for difficult, dirty and dangerous task, it also assist in human robot Interaction Tasks. The modern robot adopts AI based designs to perform in unstructured world.  Cognitive Developmental Robotics (CDR), which exists at the intersection of robotics and cognitive science, neuroscience and developmental psychology, is an innovative approach for design of cognitive robots. This approach has interdisciplinary characteristics and aims to tackle limitations that exist in current robot design. By following CDR principle, cognitive robotics concentrates on design of artificial architectures and computational models in which cognitive agents/robots higher order behavior is emerged from development and layered on former ability. The main methodologies of this approach are development learning and grounded cognition. The Consideration of these two methodologies in design of robot enables it for general tasks: Human Robot Interaction, Surgical robotics, Space exploration, Human subject centered studies (Neuroscience and psychology), Agriculture robotics

    \textbf{Keywords:} \textit{Robotics, Cognitive Science, Developmental Learning, Intrinsic Motivation, Extrinsic Motivation}

    % Acknowledgments
    \newpage
    \pdfbookmark[0]{Acknowledgments}{acknowledgments}
    \chapter*{Acknowledgments}
I would like to show my appreciations to my thesis supervisor, Dr. Nadia Rashid (Assistant Professor, Department of Computer Systems Engineering, Faculty of Engineering, The Islamia University of Bahawalpur) for his unconditional support, guidance and appreciation. It is also acknowledged that the open source databases are utilized in this research. I also acknowledge my family for their long patience and prayers.

    \newpage
    \pdfbookmark[0]{Contents}{contents}
    \tableofcontents

    \listoffigures
    \listoftables

    \newpage
    \chapter*{List of Abbreviations and Symbols}
    \section*{Abbreviations}
    \begin{acronyms}
    \item[LIS] Locked-in Syndrome
    \item[QoL] Quality of Life
    \item[BCI] Brain-computer Interface
    \item[EEG] Electroencephalography, Electroencephalographic
    \item[NIRS] Near Infrared Spectroscopy
    \end{acronyms}

    \newpage
    \pagenumbering{arabic}

    % Include all chapter files
    \include{tex/chapter1}
    \include{tex/chapter2}

    \phantomsection
    \addcontentsline{toc}{chapter}{References}
    \bibliography{bibliography}
\end{document}